\documentclass[a4paper,12pt]{article}
\usepackage[margin=2.5cm]{geometry}
\usepackage{graphicx}
\usepackage{booktabs}
\usepackage{amsmath}
\usepackage{amssymb}
\usepackage{caption}
\usepackage{subcaption}
\usepackage{hyperref}
\usepackage{xcolor}
\usepackage{tabularx}

\title{\textbf{Result 9: Comparative Performance Radar Chart}\\
\large ACLIF --- Autonomous Cross-Layer Intelligence Framework}
\author{B.~Tharuni~Sri~Sai\\
\small School of Computer Science and Engineering, VIT-AP University, Amaravati, India\\
\small ORCID: 0009-0004-9561-8985}
\date{\today}

\begin{document}
\maketitle
\tableofcontents
\newpage

\section{Experiment Overview}
This document presents the detailed analysis for \textbf{Result 9}: Comparative Performance Radar Chart,
as part of the ACLIF cross-layer IoT optimization framework evaluation.

\section{Result Figure}
\begin{figure}[htbp]
  \centering
  \includegraphics[width=0.92\textwidth]{../figures/fig_radar.png}
  \caption{Comparative Performance Radar Chart. Results obtained from NS3 v3.38 simulations over a
  100-node heterogeneous IoT network ($200\times200$~m$^2$), averaged over 10 runs
  with 95\% confidence interval shading.}
  \label{fig:radar}
\end{figure}

\section{Analysis and Discussion}
The radar chart provides a six-dimensional performance summary normalised to $[0,1]$
(higher is better for all axes, with inverse scaling applied to delay, energy,
collision, and overhead metrics). ACLIF's polygon consistently encloses the largest
area across all six dimensions, visually confirming its Pareto superiority.
FDRL and IPSO perform competitively on delay and throughput dimensions but fall short
on energy and communication overhead, while FACSO underperforms on both throughput
and collision dimensions.

\subsubsection*{Normalised Scores (ACLIF vs Best Baseline)}
\begin{tabular}{lcc}
\toprule
Metric & ACLIF & Best Baseline \\
\midrule
Delay (inv)     & 1.00 & 0.78 (IPSO) \\
Energy (inv)    & 1.00 & 0.81 (IPSO) \\
Throughput      & 1.00 & 0.83 (IPSO) \\
Collision (inv) & 1.00 & 0.76 (IPSO) \\
Comp. OH (inv)  & 0.77 & 1.00 (IPSO) \\
Comm. OH (inv)  & 0.84 & 1.00 (E-GLBR) \\
\bottomrule
\end{tabular}

\section{Experimental Configuration}
\begin{tabular}{ll}
\toprule
\textbf{Parameter} & \textbf{Value} \\
\midrule
Simulator      & NS3 v3.38 \\
Nodes          & 100 sensor nodes \\
Area           & $200 \times 200$~m$^2$ \\
PHY standard   & IEEE 802.15.4 \\
MAC protocol   & CSMA/CA \\
Packet payload & 512 bytes \\
Packet rate    & 5 pkt/s (default) \\
Runs per config & 10 (different seeds) \\
CI             & 95\% Student-$t$ \\
\bottomrule
\end{tabular}

\section{Baselines}
ACLIF is compared against: E-GLBR \cite{Benelhouri2023}, FDRL \cite{Suresh2024},
FACSO \cite{Yang2024204}, EEHCT \cite{Chaurasiya202325941}, and IPSO \cite{Zhang2025}.

\begin{thebibliography}{9}
\bibitem{Benelhouri2023} Benelhouri et al., ``E-GLBR Genetic Algorithm Routing,'' 2023.
\bibitem{Suresh2024} Suresh et al., ``FDRL Federated Deep RL Routing,'' 2024.
\bibitem{Yang2024204} Yang et al., ``FACSO Fuzzy Adaptive Cross-Layer Scheduling,'' 2024.
\bibitem{Chaurasiya202325941} Chaurasiya et al., ``EEHCT Energy-Efficient Clustering,'' 2023.
\bibitem{Zhang2025} Zhang et al., ``IPSO-Based Multi-Hop Routing,'' 2025.
\end{thebibliography}

\end{document}
